%============================== CONCLUSION ==================================
\chapter{Conclusion générale}

Ce projet de fin d'année a permis de concevoir et de développer, pour SBS Data
Factory, un \textbf{prototype fonctionnel de bout en bout} d'une plateforme web
de gestion de campagnes publicitaires \emph{drive-to-store}. Le livrable couvre
l'intégralité du parcours métier visé : connexion et navigation protégée par
rôle, tableau de bord, gestion et création de campagnes, import et validation
de la base de magasins, ciblage géographique sur carte, génération de créatives
dynamiques et de \emph{landing pages} personnalisées, reporting complet avec
export, et administration du compte --- le tout appuyé sur une API backend
\textbf{documentée et testable} (24 endpoints, Swagger interactif).

Sur le plan technique, le projet a mis en \oe uvre une \textbf{architecture
fullstack claire} (React/Vite et FastAPI), une \textbf{séparation nette des
responsabilités}, une \textbf{cartographie libre} (Leaflet/OpenStreetMap) et un
\textbf{flux de versionnement rigoureux} (branches et \emph{pull requests}).
La base de données PostgreSQL a été \textbf{préparée} (modèles et migrations)
sans bloquer le fonctionnement du prototype, et la documentation a été produite
\textbf{en continu}.

Le choix, assumé et documenté, de s'appuyer sur des \textbf{données simulées} et
une \textbf{authentification simulée} (à l'exception des brouillons de campagne
réellement persistés) a permis de \textbf{valider l'ensemble des parcours
utilisateur} avant d'investir dans le branchement d'une base réelle et d'une
sécurité de niveau production. Les limites correspondantes ont été énoncées avec
transparence, et une trajectoire d'évolution claire a été tracée.

Au-delà du résultat livré, ce stage a été une expérience formatrice sur
plusieurs plans : la \textbf{conduite itérative} d'un projet réel, la
\textbf{maîtrise d'une stack moderne} (React, FastAPI, OpenAPI, Leaflet), la
\textbf{discipline de versionnement} et de documentation, et la capacité à
faire des \textbf{choix d'ingénierie pragmatiques} en distinguant clairement ce
qui est réel de ce qui est simulé. Ces compétences, à la croisée du
développement web et du marketing digital, constituent un socle solide pour la
suite du parcours, et le prototype réalisé offre une base saine et évolutive
pour de futurs développements au sein de SBS Data Factory.
